% Part:sets-functions-relations
% Chapter: size-of-sets
% Section: introduction

\documentclass[../../../include/open-logic-section]{subfiles}

\begin{document}

\olfileid{sfr}{siz}{int}

\olsection{પ્રસ્તાવના}

૧૮૭૦ના દાયકામાં જ્યૉર્જ કૅન્ટૉરે ગણસિદ્ધાંત વિકસાવ્યો,
ત્યારે તેમના હેતુઓમાંનો એક અનંત સંગ્રહનો ખ્યાલ---મધ્યયુગના
વિચારકોના શબ્દોમાં, વાસ્તવમાં અસ્તિત્વ ધરાવતું અનંત---
સ્વીકાર્ય બનાવવાનો હતો. તેનો અગત્યનો ભાગ જુદા ગણોના
\emph{કદ}ની તેમની ચર્ચા હતી. જો $a$, $b$ અને $c$
બધા જુદા હોય, તો સહજ રીતે ગણ $\{a, b, c\}$ એ
$\{a, b\}$ કરતાં \emph{મોટો} છે. પરંતુ અનંત ગણો વિશે
શું? શું તે બધા એકબીજા જેટલા મોટા છે? એમ નથી તે જાણવા મળે છે.

અહીં પહેલો અગત્યનો ખ્યાલ પરિગણનાનો છે. કોઈ પણ સાન્ત
ગણના બધા !!{element}sની યાદી બનાવી શકીએ છીએ.
કેટલાક અનંત ગણોના પણ બધા !!{element}sની યાદી બનાવી
શકીએ, જો યાદીને જ અનંત રહેવાની છૂટ આપીએ.
આવા ગણોને !!{enumerable} કહે છે. કૅન્ટૉરનું આશ્ચર્યજનક
પરિણામ એ હતું કે કેટલાક અનંત ગણો !!{enumerable} નથી.
આ પ્રકરણના અંત સુધીમાં આપણે તે પરિણામ પૂરેપૂરું સમજીશું.

\end{document}
